\rhead{CS204: Computer Organization \& Architecture}
\phantomsection 
\section*{Computer Organization \& Architecture (CS204)}
\label{course:CS204}

\noindent \small{\hyperref[main:scheme]{$\leftarrow$ Back to Scheme}} \vspace{0.5cm}

\noindent \textbf{Credits:} 3 (3-0-0)

\subsection*{Course Content}
\noindent\textbf{Unit 1} \\
\textit{Fundamentals of Computer Organization:} Basic structure of a computer system (CPU, memory, I/O), Data representation and formats (fixed-point, floating-point, character codes), ALU operations and status flags, Register organization and register transfer language, Bus structures and types of data paths, Instruction cycle (fetch, decode, execute) and micro-operations, Introduction to instruction set design and basic types of instructions (data movement, arithmetic/logic, control).

\vspace{0.3cm}

\noindent\textbf{Unit 2} \\
\textit{8085 and 8086 Architectures:} Overview of 8-bit vs. 16-bit microprocessors, Internal architecture of 8085 (registers, ALU, control unit, buses) as a baseline model, 8086 architecture: Bus Interface Unit (BIU) and Execution Unit (EU), General-purpose, segment, index, and pointer registers, Memory segmentation and logical vs. physical addresses, Clock, reset, and basic timing concepts, Concept of a simple instruction pipeline via 8086 prefetch queue.

\vspace{0.3cm}

\noindent\textbf{Unit 3} \\
\textit{Instruction Sets, Addressing Modes, and Control:} 8086 instruction format and classification (data transfer, arithmetic, logic, control flow, string operations), Addressing modes (immediate, register, direct, register indirect, indexed, based indexed, relative), Stack organization and procedure calls/returns, Condition codes and branching, Hardwired vs. microprogrammed control units, Micro-instructions and micro-operations, Introduction to interrupt structure (maskable vs. non-maskable) and basic exception handling model.

\vspace{0.3cm}

\noindent\textbf{Unit 4} \\
\textit{Memory Hierarchy and I/O Organization:} Main memory organization (RAM, ROM, memory maps), Memory interleaving and bandwidth considerations, Cache memory hierarchy (L1/L2/L3), cache mapping techniques (direct-mapped, associative, set-associative), cache write policies and performance trade-offs, Concept of virtual memory and address translation, I/O-mapped vs. memory-mapped I/O, Programmed I/O, interrupt-driven I/O, and Direct Memory Access (DMA) for high-speed peripherals.

\vspace{0.3cm}

\noindent\textbf{Unit 5} \\
\textit{Modern Microarchitectures and Multi-Core Systems:} RISC vs. CISC philosophies in contemporary processors (x86 vs. ARM), Pipelining hazards (structural, data, control) and high-level techniques to mitigate them, Overview of superscalar execution and out-of-order pipelines, Multi-core and many-core processors, shared memory and coherence concepts, System-on-Chip (SoC) integration in modern designs, Case studies of current CPU families (Intel/AMD x86 and Apple M-series) focusing on cores, caches, integrated accelerators (GPU/Neural engines), and unified memory architectures.

\newpage
\rhead{Curriculum Schema}
